% =====================================================================
%  Proyecto Aprendizaje por Refuerzo - Entrega 1
%  Definición de estados, acciones y función de recompensa R(s,a,s')
%  Compilar con pdfLaTeX:  latexmk -pdf entrega1_definicion_mdp.tex
% =====================================================================
\documentclass[11pt,letterpaper]{article}

% ---------- Idioma y codificación ----------
\usepackage[T1]{fontenc}
\usepackage[utf8]{inputenc}
\usepackage{babel}
\babelprovide[import,main]{spanish}   % funciona aunque no esté instalado spanish.ldf
\usepackage{lmodern}

% ---------- Página ----------
\usepackage[margin=2.3cm]{geometry}
\usepackage{fancyhdr}
\usepackage{titlesec}
\usepackage{parskip}

% ---------- Matemáticas y tablas ----------
\usepackage{amsmath,amssymb}
\usepackage{booktabs}
\usepackage{tabularx}
\usepackage{array}
\usepackage{multirow}
\usepackage[table]{xcolor}
\usepackage{enumitem}
\usepackage{caption}
\usepackage[colorlinks=true,linkcolor=azul,urlcolor=azul]{hyperref}

% ---------- Colores y estilo ----------
\definecolor{azul}{HTML}{1F4E79}
\definecolor{gris}{HTML}{F2F4F7}
\titleformat{\section}{\Large\bfseries\color{azul}}{\thesection.}{0.6em}{}
\titleformat{\subsection}{\large\bfseries\color{azul}}{\thesubsection}{0.6em}{}
\captionsetup{font=small,labelfont=bf}
\setlist{itemsep=2pt,topsep=4pt}
\newcolumntype{L}{>{\raggedright\arraybackslash}X}
\newcolumntype{C}[1]{>{\centering\arraybackslash}p{#1}}
\renewcommand{\arraystretch}{1.25}

\pagestyle{fancy}
\fancyhf{}
\lhead{\small Aprendizaje por Refuerzo -- Proyecto, Entrega 1}
\rhead{\small Consolidación de picking en un CD}
\cfoot{\small \thepage}

% ---------- Macros ----------
\newcommand{\LOTE}[1]{\mathsf{LOTE}_{#1}}
\newcommand{\URG}{\mathsf{URGENTE}}
\newcommand{\ESP}{\mathsf{ESPERAR}}
\newcommand{\sfin}{s_{\mathrm{fin}}}
\newcommand{\umax}{u_{\max}}

% =====================================================================
\begin{document}

\begin{center}
  {\small Universidad de los Andes \textperiodcentered{} Maestría en Inteligencia Artificial}\\[2pt]
  {\small Aprendizaje por Refuerzo \textperiodcentered{} Prof. Nicolás Cardozo}\\[10pt]
  {\LARGE\bfseries\color{azul} Agente de RL para la consolidación de órdenes\\[2pt] de picking en un centro de distribución}\\[8pt]
  {\large Entrega 1: estados, acciones y función de recompensa}\\[8pt]
  Miguel Ángel Muñoz Layton \quad\textperiodcentered\quad Daniel F. Barrera Armendaris\\[2pt]
  {\small Noviembre de 2026}
\end{center}

\vspace{4pt}
\hrule
\vspace{8pt}

% =====================================================================
\section{Descripción de la aplicación}

Un agente de aprendizaje por refuerzo decide, cada vez que un operario queda libre,
qué trabajo de picking liberar en un centro de distribución (CD): un lote consolidado de
una zona de racks, una orden urgente individual, o esperar a que lleguen más órdenes.
El objetivo es reducir los recorridos a los racks agrupando ítems repetidos entre órdenes,
sin incumplir el deadline de ninguna orden.

\textbf{El problema.} A lo largo del turno llegan órdenes de compra. Cada orden tiene líneas
(SKU y cantidad), cada SKU tiene una ubicación fija en un rack y cada orden tiene un deadline
(hora de corte del transportista). Si cada orden se recoge por separado y 5 órdenes piden el
mismo SKU, un operario va 5 veces a la misma ubicación. El agente agrupa esas líneas para hacer
una sola visita y recoger las 5 unidades.

\textbf{La tensión que aprende el agente.} Esperar a que lleguen más órdenes aumenta la
consolidación, pero consume holgura frente al deadline y deja operarios ociosos. Una regla fija
no resuelve bien ese balance cuando la demanda, el WIP y la cantidad de operarios cambian durante
el día; por eso se plantea como un problema de RL.

\textbf{Supuestos del entorno.}
\begin{itemize}
  \item Los racks están organizados en 4 zonas, $Z=\{A,B,C,D\}$; cada zona es un grupo de pasillos contiguos.
  \item Un episodio es un turno de 8 horas.
  \item Cada operario lleva un carro con capacidad de $C_{\max}=20$ líneas.
  \item El número de operarios puede variar en el turno (pausas, ausencias); el agente observa cuántos están libres.
  \item Las órdenes llegan de forma estocástica, con picos de demanda durante el turno; el agente no conoce las órdenes futuras.
  \item Pocos SKU concentran la mayor parte de la demanda (curva ABC), por lo que varias órdenes suelen pedir los mismos SKU.
\end{itemize}

\textbf{Lote consolidado.} Al liberar un lote en la zona $z$, el sistema toma las líneas pendientes
de $z$ ordenadas por deadline más cercano, hasta 20 líneas, y las agrupa por ubicación: todas las
unidades del mismo SKU se recogen en una sola visita. Después se separan por orden en una mesa de
consolidación (\emph{put-wall}) antes del despacho.

\textbf{Momentos de decisión.} Hay una decisión (a) cada vez que un operario queda libre y (b) cada
$\Delta t = 2$ min en cualquier caso: el reloj nunca avanza más de $\Delta t$ por decisión. Si todos los
operarios están ocupados, la única acción posible es esperar. Cada acción asigna como máximo un
operario; si quedan más operarios libres y hay trabajo, se toma otra decisión de inmediato sin que avance el reloj.

\subsection{Comportamiento deseado y no deseado del agente}

\begin{table}[h]
\centering\small
\begin{tabularx}{\textwidth}{@{}p{0.7cm}Lp{2.1cm}L@{}}
\toprule
\textbf{\#} & \textbf{Comportamiento} & \textbf{Tipo} & \textbf{Cómo lo controla la recompensa} \\
\midrule
D1 & Recoger en una sola visita un SKU que piden varias órdenes & Deseado & $+2$ por cada visita ahorrada (líneas del lote menos ubicaciones visitadas) \\
D2 & Atender primero las órdenes con holgura crítica ($\le 15$ min) & Deseado & $-1$ por zona crítica cada vez que avanza el reloj; $-20$ por orden vencida \\
D3 & Mantener ocupados a los operarios cuando hay trabajo & Deseado & $\ESP$ con operarios libres cuesta $-1$ o $-3$, más el costo de WIP \\
D4 & Esperar brevemente para formar lotes más grandes si llegan muchas órdenes y hay holgura & Deseado & $\ESP$ no se penaliza (0) si $\lambda=2$ y la consolidación aún es baja; las visitas ahorradas del lote siguiente pagan la espera \\
D5 & Mantener controlado el WIP & Deseado & $-0{,}05$ por cada línea pendiente sin asignar, por minuto \\
D6 & Despachar todas las órdenes antes del fin de turno & Deseado & $+3$ por orden a tiempo; $-30$ por orden pendiente al cierre \\
\midrule
N1 & Ir varias veces a la misma ubicación por órdenes distintas & No deseado & $\URG$ sin urgencia cuesta $-4$ extra; cada viaje cuesta $-1$ y uno sin SKU repetidos no ahorra visitas \\
N2 & Retener órdenes esperando consolidar hasta que vencen & No deseado & $-20$ por orden vencida, más que lo que ahorra consolidar todas sus líneas ($\le +12$) \\
N3 & Dejar operarios ociosos con órdenes críticas pendientes & No deseado & $\ESP$ con una zona crítica cuesta $-10$ \\
N4 & Olvidar zonas con poca coincidencia de SKU (inanición) & No deseado & Su urgencia $u_z$ sube: $-1$ por avance mientras sea crítica y $-10$ a los lotes de otras zonas \\
N5 & Atender una zona cómoda mientras otra tiene una orden crítica & No deseado & $-10$ extra a $\LOTE{z}$ si otra zona es crítica y $z$ no \\
N6 & Asignar trabajo sin operarios libres o a una zona sin líneas & No deseado & Acción enmascarada; si se ejecuta, $-50$ \\
\bottomrule
\end{tabularx}
\caption{Comportamientos del agente y el componente de recompensa que los incentiva o castiga.}
\end{table}

% =====================================================================
\section{Conjunto de estados}

El estado es un vector de 14 variables discretas: tres por zona (WIP, consolidación y urgencia),
más los operarios libres y la tasa de llegada de órdenes.
\[
  s = (w_A, w_B, w_C, w_D,\; c_A, c_B, c_C, c_D,\; u_A, u_B, u_C, u_D,\; o,\; \lambda),
  \qquad z \in Z = \{A,B,C,D\}.
\]
Para abreviar, en el resto del documento se escribe
\[ s = (w_A,w_B,w_C,w_D \mid c_A,c_B,c_C,c_D \mid u_A,u_B,u_C,u_D \mid o,\lambda). \]

\begin{table}[h]
\centering\small
\begin{tabularx}{\textwidth}{@{}p{1cm}Lp{1.9cm}L@{}}
\toprule
\textbf{Var.} & \textbf{Significado} & \textbf{Dominio} & \textbf{Discretización} \\
\midrule
$w_z$ & WIP de la zona $z$: líneas pendientes aún no asignadas a un operario & $\{0,1,2,3\}$ &
  0 = 0 líneas; 1 = 1--5; 2 = 6--15; 3 = 16 o más (zona saturada) \\
$c_z$ & Potencial de consolidación: máximo número de órdenes pendientes que piden un mismo SKU en $z$ & $\{0,1,2\}$ &
  0 = ningún SKU repetido; 1 = 2--4 órdenes comparten un SKU; 2 = 5 o más \\
$u_z$ & Urgencia: holgura mínima entre las órdenes con líneas pendientes en $z$ (holgura = deadline $-$ hora actual $-$ tiempo estimado de recogida) & $\{0,1,2,3\}$ &
  0 = zona sin líneas; 1 = más de 60 min; 2 = 15--60 min; 3 = 15 min o menos, incluye vencidas (crítica) \\
$o$ & Operarios libres en el momento de la decisión & $\{0,1,2,3\}$ & 0, 1, 2; 3 = tres o más \\
$\lambda$ & Órdenes llegadas en los últimos 15 min & $\{0,1,2\}$ & 0 = menos de 5; 1 = 5--15; 2 = más de 15 \\
\bottomrule
\end{tabularx}
\caption{Variables del estado y sus dominios.}
\end{table}

\textbf{Restricción de consistencia.} Una zona sin WIP no tiene consolidación ni urgencia:
\[
  w_z = 0 \iff c_z = 0 \ \wedge\ u_z = 0 .
\]
Por tanto, cada zona tiene $1 + 3\cdot3\cdot3 = 28$ combinaciones válidas de $(w_z,c_z,u_z)$.

\textbf{Estado terminal.} $\sfin$ representa el fin del turno (8 h). Ahí termina el episodio y no hay acciones.

\textbf{Tamaño del espacio.}
\[
  |S| = 28^4 \cdot 4 \cdot 3 + 1 = 7\,375\,873 \text{ estados.}
\]
Es finito pero grande para una tabla $Q$ completa. Una primera versión tabular con 2 zonas
tiene $28^2\cdot 12 + 1 = 9\,409$ estados.

El estado abstrae la posición exacta de cada SKU y la identidad de cada orden; esa información
la usa el simulador para armar los lotes y contar eventos, no el agente para decidir.

\subsection{Clases de estados relevantes}

Para definir las acciones y las recompensas basta distinguir las siguientes clases. Sea
$Z^{+}(s) = \{z \in Z : w_z \ge 1\}$ el conjunto de zonas con trabajo y
$\umax = \max_{z} u_z$.

\begin{table}[h]
\centering\small
\begin{tabularx}{\textwidth}{@{}p{1cm}p{5.2cm}L@{}}
\toprule
\textbf{Clase} & \textbf{Condición} & \textbf{Interpretación} \\
\midrule
$S_0$ & $s = \sfin$ & Fin del turno \\
$S_1$ & $o = 0$ & Todos los operarios ocupados \\
$S_2$ & $o \ge 1$ y $Z^{+}(s)=\emptyset$ & Hay operarios libres pero no hay trabajo pendiente \\
$S_3$ & $o \ge 1$, $Z^{+}(s)\neq\emptyset$, $\umax = 3$ & Hay al menos una orden crítica \\
$S_4$ & $o \ge 1$, $Z^{+}(s)\neq\emptyset$, $\umax \le 2$ & Hay trabajo sin urgencia crítica \\
\bottomrule
\end{tabularx}
\caption{Partición del espacio de estados usada en acciones y recompensas.}
\end{table}

% =====================================================================
\section{Conjunto de acciones}

El agente tiene 6 acciones: un lote consolidado por cada zona, recoger una orden urgente completa,
o esperar.
\[
  A = \{\LOTE{A}, \LOTE{B}, \LOTE{C}, \LOTE{D}, \URG, \ESP\}.
\]

\begin{table}[h]
\centering\small
\begin{tabularx}{\textwidth}{@{}p{2.4cm}Lp{4.2cm}@{}}
\toprule
\textbf{Acción} & \textbf{Qué hace} & \textbf{Aplicable si} \\
\midrule
$\LOTE{z}$, $z\in Z$ & Asigna un operario libre a un lote de la zona $z$: hasta 20 líneas por deadline más cercano, agrupadas por ubicación (una visita por SKU) & $o\ge 1$ y $w_z \ge 1$ \\
$\URG$ & Asigna un operario libre a recoger completa la orden con menor holgura de todo el CD, en todas sus zonas, sin mezclarla con otras órdenes & $o\ge1$ y $Z^{+}(s)\neq\emptyset$ \\
$\ESP$ & No asigna trabajo; el reloj avanza hasta el siguiente evento (máximo $\Delta t = 2$ min) y pueden llegar órdenes nuevas & Siempre, salvo en $\sfin$ \\
\bottomrule
\end{tabularx}
\caption{Acciones del agente.}
\end{table}

\textbf{Acciones aplicables por estado.} Formalmente,
\[
A(s) =
\begin{cases}
  \emptyset & s \in S_0 \\[2pt]
  \{\ESP\} & s \in S_1 \cup S_2 \\[2pt]
  \{\LOTE{z} : z \in Z^{+}(s)\} \cup \{\URG,\ \ESP\} & s \in S_3 \cup S_4
\end{cases}
\]
Así, en $S_3 \cup S_4$ el agente tiene entre 3 y 6 acciones disponibles (según cuántas zonas
tengan WIP). Las acciones fuera de $A(s)$ se enmascaran; si una implementación no las enmascara,
ejecutar una acción inválida da $-50$ y el estado no cambia ($s'=s$).

\subsection{Dinámica de transición: cómo se llega a \texorpdfstring{$s'$}{s'}}

La recompensa se define sobre la tripla $(s,a,s')$, así que conviene precisar qué cambia de $s$
a $s'$ con cada acción. Sea $\tau(s,a,s') \in \{0,1\}$ un indicador de que el reloj avanzó
entre $s$ y $s'$.

\begin{table}[h]
\centering\small
\begin{tabularx}{\textwidth}{@{}p{2.1cm}LL@{}}
\toprule
\textbf{Acción} & \textbf{Efecto inmediato} & \textbf{Avance del reloj} \\
\midrule
$\LOTE{z}$ & $o \to o-1$; salen de $z$ hasta 20 líneas, y $w_z, c_z, u_z$ se recalculan con las líneas que quedan & $\tau=0$ si $o-1\ge1$ y aún hay WIP (nueva decisión inmediata); si no, $\tau=1$ y el reloj avanza $\Delta t$ o hasta que un operario quede libre, lo que ocurra primero \\
$\URG$ & $o \to o-1$; salen las líneas de la orden más urgente en todas las zonas donde tenga líneas, y esas zonas se recalculan & Igual que $\LOTE{z}$ \\
$\ESP$ & Ninguno & $\tau=1$: el reloj avanza $\Delta t$ o hasta que se libere un operario, lo que ocurra primero \\
\bottomrule
\end{tabularx}
\caption{Transiciones por acción.}
\end{table}

Así, ninguna transición dura más de $\Delta t$ y todas acumulan costos de intervalos comparables.
Cuando $\tau = 1$ ocurren los eventos estocásticos del ambiente: llegan órdenes ($w$, $c$ y
$\lambda$ pueden subir), la holgura de las órdenes pendientes disminuye ($u$ puede subir),
operarios terminan su lote ($o$ puede subir) y órdenes se completan o vencen. Si en el avance se
alcanzan las 8 h, $s'=\sfin$.

% =====================================================================
\section{Función de recompensa \texorpdfstring{$R(s,a,s')$}{R(s,a,s')}}

La recompensa de cada tripla suma una parte inmediata, que depende del estado de salida y de la
acción, y una parte de transición, que depende de lo que ocurrió entre $s$ y $s'$:
\[
  R(s,a,s') = R_{\mathrm{acc}}(s,a) + R_{\mathrm{tr}}(s,a,s').
\]
La consolidación no se paga con un bono fijo por lote, sino por las \emph{visitas ahorradas}: si un lote
tiene $\ell$ líneas y visita $v$ ubicaciones distintas, ahorra $n_{\mathrm{ahorr}}=\ell-v$ visitas. Así,
despachar muchos lotes pequeños no da más recompensa que uno grande, y esperar a que se acumulen SKU
repetidos sí puede pagar.

\subsection{Recompensa inmediata \texorpdfstring{$R_{\mathrm{acc}}(s,a)$}{R\_acc(s,a)}}

Cada viaje despachado ($\LOTE{z}$ o $\URG$) cuesta $-1$. Las reglas de $\ESP$ se evalúan en orden; aplica la primera que se cumpla.

\begin{table}[h]
\centering\small
\begin{tabularx}{\textwidth}{@{}p{2cm}Lc@{}}
\toprule
\textbf{Acción} & \textbf{Condición sobre $s$} & \textbf{$R_{\mathrm{acc}}(s,a)$} \\
\midrule
\multirow{2}{*}{$\LOTE{z}$} & Ninguna otra zona es crítica, o $z$ también lo es & $-1$ \\
 & Otra zona es crítica ($\exists z'\neq z: u_{z'}=3$) y $z$ no ($u_z<3$): descuido & $-1-10=-11$ \\
\midrule
\multirow{2}{*}{$\URG$} & $\umax \ge 2$ & $-1$ \\
 & $\umax = 1$ (no hay urgencia: desperdicia consolidación) & $-1-4=-5$ \\
\midrule
\multirow{6}{*}{$\ESP$} & $s\in S_1$: $o = 0$ (no hay alternativa) & $0$ \\
 & $s\in S_2$: no hay trabajo pendiente & $0$ \\
 & $s\in S_3$: operario ocioso con una orden crítica & $-10$ \\
 & $s\in S_4$, $\lambda = 2$ y $c_z \le 1\ \forall z$ (esperar puede formar lotes más consolidados) & $0$ \\
 & $s\in S_4$ y $\lambda = 0$ (casi no llegan órdenes, esperar no sirve) & $-3$ \\
 & $s\in S_4$ en cualquier otro caso (operario ocioso) & $-1$ \\
\bottomrule
\end{tabularx}
\caption{Recompensa inmediata por acción.}
\end{table}

Ignorar una orden crítica cuesta $-10$ tanto si el operario espera como si se le asigna otra zona.

\subsection{Recompensa de transición \texorpdfstring{$R_{\mathrm{tr}}(s,a,s')$}{R\_tr(s,a,s')}}

Sean $n_{\mathrm{ahorr}}$ las visitas ahorradas por el lote asignado (0 si $a=\ESP$); $n_{\mathrm{ok}}$ y
$n_{\mathrm{venc}}$ las órdenes completadas a tiempo y las que vencieron sin completarse durante $s \to s'$;
$A_{\mathrm{wip}}$ el área de WIP del intervalo, es decir, la suma sobre el intervalo de las líneas pendientes
sin asignar por los minutos que permanecen así ($\le \Delta t$ minutos, porque el reloj nunca avanza más); y
$n_{\mathrm{pend}}$ las órdenes sin despachar al llegar a $\sfin$. Como el estado agregado no guarda la
identidad de cada orden, el simulador registra estos conteos en cada transición; para una tripla concreta su
valor es conocido.
\[
\begin{aligned}
  R_{\mathrm{tr}}(s,a,s') = {}& 2\,n_{\mathrm{ahorr}} + 3\,n_{\mathrm{ok}} - 20\,n_{\mathrm{venc}} - 0{,}05\,A_{\mathrm{wip}} \\
   & - \tau(s,a,s')\,\mathbf{1}[s'\neq\sfin]\sum_{z\in Z}\mathbf{1}[u'_z = 3]
   - 30\,n_{\mathrm{pend}}\,\mathbf{1}[s'=\sfin].
\end{aligned}
\]

\begin{table}[h]
\centering\small
\begin{tabularx}{\textwidth}{@{}Lc@{}}
\toprule
\textbf{Evento en la transición $(s,a,s')$} & \textbf{Recompensa} \\
\midrule
Cada visita ahorrada por el lote asignado ($\ell - v$) & $+2$ \\
Cada orden completada antes de su deadline & $+3$ \\
Cada orden que vence sin completarse (se cuenta una sola vez por orden) & $-20$ \\
Cada línea pendiente sin asignar, por minuto transcurrido (costo de WIP) & $-0{,}05$ \\
Cada zona con $u'_z = 3$ en $s'$, solo si el reloj avanzó ($\tau=1$) y $s'\neq\sfin$ & $-1$ \\
Llegada a $s'=\sfin$: cada orden pendiente sin despachar & $-30$ \\
Acción inválida ($a\notin A(s)$, si no hay enmascaramiento); en este caso $s'=s$ y no aplica nada más & $-50$ \\
\bottomrule
\end{tabularx}
\caption{Recompensas de transición.}
\end{table}

\textbf{Justificación de los valores.}
\begin{itemize}
  \item Vencer una orden ($-20$) cuesta más de lo que se gana consolidando sus líneas (una orden tiene a lo
    sumo 6 líneas: $\le +12$), así que el agente no debería sacrificar un deadline por consolidar.
  \item El costo de WIP hace que esperar nunca sea gratis: con 20 líneas pendientes, esperar $\Delta t=2$ min
    cuesta $-2$, lo mismo que una visita ahorrada. Esperar solo conviene si las órdenes que llegan
    comparten SKU con las pendientes, que es justamente el balance que el agente debe descubrir.
  \item Cobrar por visita ahorrada, y no por lote, evita que el agente gane recompensa fragmentando el trabajo
    en muchos lotes pequeños; el $-1$ por viaje refuerza lo mismo.
\end{itemize}

\textbf{Descuento.} Como el reloj avanza a lo sumo $\Delta t$ por decisión pero algunas decisiones son
inmediatas ($\tau=0$), el descuento se aplica por tiempo transcurrido y no por decisión:
$\gamma^{\,\Delta_{\mathrm{min}}/\Delta t}$ con $\gamma = 0{,}95$ por cada $\Delta t = 2$ min. Una
decisión inmediata no se descuenta, de modo que asignar varios operarios en el mismo instante no aplaza
artificialmente las recompensas futuras.

\subsection{Triplas \texorpdfstring{$(s,a,s')$}{(s,a,s')} de ejemplo}

La tabla aplica la función completa a transiciones concretas que cubren todas las clases de estado
y todas las acciones. Notación: $s=(w_A,w_B,w_C,w_D \mid c_A,c_B,c_C,c_D \mid u_A,u_B,u_C,u_D \mid o,\lambda)$.
Para calcular el costo de WIP se indica cuántas líneas hay pendientes; las órdenes que llegan en el intervalo
se suponen al final de este.

\begin{table}[h]
\centering\footnotesize
\setlength{\tabcolsep}{4pt}
\begin{tabularx}{\textwidth}{@{}p{0.5cm}p{5.0cm}p{1.6cm}Lc@{}}
\toprule
\textbf{\#} & \textbf{Estado de salida $s$ y de llegada $s'$} & \textbf{Acción $a$} & \textbf{Cálculo} & \textbf{$R(s,a,s')$} \\
\midrule
T1 & $s=(1,2,0,0|0,2,0,0|3,1,0,0|2,1)$\newline $s'=(0,2,0,0|0,2,0,0|0,1,0,0|1,1)$ & $\LOTE{A}$ & A tiene 3 líneas críticas de 3 SKU distintos: viaje $-1$, sin visitas ahorradas; queda un operario libre y hay WIP en B, $\tau=0$ & $-1$ \\
\addlinespace[2pt]
T2 & $s=(0,2,0,0|0,2,0,0|0,1,0,0|1,1)$\newline $s'=(0,0,1,0|0,0,0,0|0,0,1,0|1,1)$ & $\LOTE{B}$ & B tiene 8 líneas: 5 órdenes piden el mismo SKU más 3 SKU distintos, 4 visitas: viaje $-1$, ahorro $4\times2=+8$; en los 2 min termina el lote de A y se completa su orden ($+3$); no quedan líneas sin asignar; llegan órdenes a C & $+10$ \\
\addlinespace[2pt]
T3 & $s=(2,1,0,0|1,0,0,0|1,1,0,0|1,1)$\newline $s'=(2,0,0,0|1,0,0,0|1,0,0,0|0,1)$ & $\URG$ & $\umax=1$: $-1-4$; la orden tenía 2 líneas en A y 1 en B; quedan 6 líneas en A: WIP $6\times2\times0{,}05=-0{,}6$ & $-5{,}6$ \\
\addlinespace[2pt]
T4 & $s=(0,0,1,0|0,0,0,0|0,0,3,0|1,1)$\newline $s'=(0,0,1,0|0,0,0,0|0,0,3,0|1,1)$ & $\ESP$ & $s\in S_3$: $-10$; la orden crítica de C (1 línea) vence: $-20$; C sigue crítica: $-1$; WIP $-0{,}1$ & $-31{,}1$ \\
\addlinespace[2pt]
T5 & $s=(2,1,0,0|1,1,0,0|1,1,0,0|1,2)$\newline $s'=(2,1,0,0|2,1,0,0|1,1,0,0|1,2)$ & $\ESP$ & $s\in S_4$, $\lambda=2$, $c_z\le1$: $0$; 11 líneas pendientes: WIP $-1{,}1$; llegan órdenes y $c_A$ sube a 2 & $-1{,}1$ \\
\addlinespace[2pt]
T6 & $s=(3,1,0,0|1,0,0,0|2,1,0,0|0,1)$\newline $s'=(3,1,0,0|1,0,0,0|2,1,0,0|1,1)$ & $\ESP$ & $s\in S_1$: $0$; se completa una orden: $+3$; 22 líneas pendientes: WIP $-2{,}2$ & $+0{,}8$ \\
\addlinespace[2pt]
T7 & $s=(0,2,1,0|0,2,0,0|0,1,3,0|1,1)$\newline $s'=(0,0,1,0|0,0,0,0|0,0,3,0|0,1)$ & $\LOTE{B}$ & Lote de T2 ($-1+8$); descuido (C es crítica): $-10$; vence la orden de C: $-20$; C sigue crítica: $-1$; WIP $-0{,}1$ & $-24{,}1$ \\
\addlinespace[2pt]
T8 & $s=(1,0,0,0|0,0,0,0|3,0,0,0|0,0)$\newline $s'=\sfin$ & $\ESP$ & $s\in S_1$: $0$; WIP $2\times2\times0{,}05=-0{,}2$; termina el turno con 2 órdenes pendientes: $2\times(-30)$ & $-60{,}2$ \\
\addlinespace[2pt]
T9 & $s=(1,2,0,0|0,2,0,0|3,1,0,0|2,1)$\newline $s'=s$ (sin cambio) & $\LOTE{C}$ & $w_C=0$: acción inválida & $-50$ \\
\bottomrule
\end{tabularx}
\caption{Ejemplos de recompensa para triplas $(s,a,s')$.}
\end{table}

Las triplas T1 y T2 muestran el caso que motiva la aplicación: con una orden crítica en A y 5 órdenes
que piden el mismo SKU en B, atender primero la urgencia cuesta $-1$, mientras que ir a B de inmediato
daría $-1-10+8=-3$ y arriesgaría el vencimiento de A. Luego el agente recoge el SKU repetido en una sola
visita (T2, $+8$ por ahorro), en vez de ir 5 veces a la misma ubicación. T7 muestra el costo de ignorar una
urgencia por perseguir consolidación.

\end{document}
